


\section{Distributional theory and uncertainty quantification}
\label{sec:distribution-theory}


Thus far, we have demonstrated intriguing statistical performance of estimators developed based on spectral methods. 
As one can anticipate, the quality of a spectral estimator is largely affected by the imperfectness of data generating mechanisms (e.g., noise corruption, missing data). 
The uncertainty of the estimator due to these factors would inevitably influence any subsequent decision making based on it. 
Viewed in this light, it is recommended to accompany the estimator in hand with valid measures of uncertainty (or ``confidence''), 
in order to better inform decision makers.  


Take the low-rank matrix estimation problem in Section~\ref{sec:setup-general-theory} for instance:  
an important uncertainty quantification task can be posed as the construction of a valid confidence
interval---based on the spectral estimator---that is likely to cover an unseen entry of the matrix of interest $\bm{M}^{\star}$. 
More precisely, for any location $(i,j)$ and any target coverage level $1-\alpha \in (0,1)$ (e.g., 95\%), we aim to identify a short interval---denoted by $\mathsf{CI}_{i,j}^{1-\alpha}$---based on the spectral estimator such that
%
\begin{equation}
	\mathbb{P}\big( M_{i,j}^{\star} \in \mathsf{CI}_{i,j}^{1-\alpha} \big) \approx 1-\alpha,  
\end{equation}
%
which essentially augments a point estimate into an interval that is guaranteed to cover the unknown with the pre-specified target probability.  
Note that the problem of constructing a valid confidence interval falls within the realm of {\em statistical inference} in the statistics literature,
which constitutes an important step beyond statistical estimation. 
Accomplishing this task in high dimension often calls for a refined statistical reasoning toolbox that offers quantitative distributional characterizations of the estimator. 

%The $\ell_{\infty}$ and $\ell_{2,\infty}$ theory presented in this chapter becomes particularly effective for this purpose.    


%calls for a modern suite
%of statistical reasoning tools, with the aim of providing comprehensive uncertainty evaluations for the estimation
%algorithms in use.
%i
%


%quantitative characterization of 

%statistical validity 



\subsection{Entrywise distributional guarantees}
\label{sec:distribution-general}


As a natural starting point to build confidence intervals, 
we seek to develop comprehensive understanding about the distribution of the spectral estimator. 
In general, obtaining a non-asymptotic yet tractable distributional characterization of a nonconvex estimator like the spectral method 
could be remarkably challenging. 
Fortunately, the $\ell_{\infty}$ and $\ell_{2,\infty}$ perturbation theory introduced previously (e.g., Theorem~\ref{thm:UsgnH-Ustar-MUstar-general}) allows one to make progress for some important scenarios. 



Let us revisit the setting in Section~\ref{sec:setup-general-theory}, and consider the following estimator of the unknown low-rank matrix $\bm{M}^{\star}$:
%
\begin{align}
	\widehat{\bm{M}}=\big[\widehat{M}_{i,j} \big]_{1\leq i,j\leq n}=\bm{U}\bm{\Lambda}\bm{U}^{\top} ,
	\label{eq:estimator-M-hat-distribution}
\end{align}
%
obtained via the spectral method. The aim is to develop tractable distributional guarantees for each entry of $\widehat{\bm{M}}-\bm{M}^{\star}$. 


Towards this end, we first examine whether our previous results shed light on certain distributional properties of $\widehat{\bm{M}}-\bm{M}^{\star}$. 
Informally,  Theorem~\ref{thm:UsgnH-Ustar-MUstar-general} (in particular, \eqref{eq:UsgnH-MUstar-bound-theorem-general}) reveals that
%
\begin{align}
	\bm{U}\mathsf{sgn}(\bm{H}) \approx \bm{M} \bm{U}^{\star} \big(\bm{\Lambda}^{\star}\big)^{-1} 
	= \bm{U}^{\star} + \bm{E} \bm{U}^{\star} \big(\bm{\Lambda}^{\star}\big)^{-1}. 
	\label{eq:U-sgn-H-first-order-approx-discuss}
\end{align}
%
Assuming tightness of this first-order approximation, one further derives 
%
\begin{align}
\bm{U}\bm{\Lambda}\bm{U}^{\top}-\bm{M}^{\star} & \overset{(\mathrm{i})}{\approx}\bm{U}\mathsf{sgn}(\bm{H})\bm{\Lambda}^{\star}\big(\bm{U}\mathsf{sgn}(\bm{H})\big)^{\top}-\bm{U}^{\star}\bm{\Lambda}^{\star}\bm{U}^{\star\top} \notag\\
 & \overset{(\mathrm{ii})}{\approx}\big(\bm{U}\mathsf{sgn}(\bm{H})-\bm{U}^{\star}\big)\bm{\Lambda}^{\star}\bm{U}^{\star\top}+\bm{U}^{\star}\bm{\Lambda}^{\star}\big(\bm{U}\mathsf{sgn}(\bm{H})-\bm{U}^{\star}\big)^{\top} 
 \notag\\
 & \overset{(\mathrm{iii})}{\approx}\bm{E}\bm{U}^{\star}\big(\bm{\Lambda}^{\star}\big)^{-1}\bm{\Lambda}^{\star}\bm{U}^{\star\top}+\bm{U}^{\star}\bm{\Lambda}^{\star}\big(\bm{E}\bm{U}^{\star}\big(\bm{\Lambda}^{\star}\big)^{-1}\big)^{\top} \notag\\
 & =\bm{E}\bm{U}^{\star}\bm{U}^{\star\top}+\bm{U}^{\star}\bm{U}^{\star\top}\bm{E}, 
	\label{eq:approx-ULambdaU-Mstar-discuss}
\end{align}
%
where (i) holds as long as $\mathsf{sgn}(\bm{H})\bm{\Lambda}^{\star} \mathsf{sgn}(\bm{H})^{\top} \approx \bm{\Lambda}$ (which has already been illuminated in the analysis 
of Corollary~\ref{cor:entrywise-error-general} and will be solidified momentarily),  
(ii) is obtained by dropping the higher-order term $\big(\bm{U}\mathsf{sgn}(\bm{H})-\bm{U}^{\star}\big)\bm{\Lambda}^{\star}\big(\bm{U}\mathsf{sgn}(\bm{H})-\bm{U}^{\star}\big)^{\top}$, and (iii) relies upon the approximation \eqref{eq:U-sgn-H-first-order-approx-discuss}.  


Given that \eqref{eq:approx-ULambdaU-Mstar-discuss} is a linear map of the noise matrix $\bm{E}$, 
this essentially forms a first-order approximation of $\widehat{\bm{M}}$,
which in turn enables a tractable distributional theory for $\widehat{\bm{M}}$. 
Observe that each entry of the matrix in \eqref{eq:approx-ULambdaU-Mstar-discuss} 
is a weighted superposition of the independent zero-mean entries of $\bm{E}$.   
Equipped with this observation, some variant of the central limit theorem
suggests that each entry of $\widehat{\bm{M}} - \bm{M}^{\star}$ is approximately zero-mean Gaussian,
as formalized by the theorem below. For notational convenience, we shall define a projection matrix
	%
	\begin{equation}
	%\bm{P}\coloneqq\bm{U}^{\star}\bm{U}^{\star\top},
		\bm{P}^{\star}= \big[ P_{i,j}^{\star} \big]_{1\leq i,j\leq n} \coloneqq \bm{U}^{\star}\bm{U}^{\star\top},
		\label{eq:defn-P-UU-T}
	\end{equation}	
	%
and impose a lower bound requirement on the noise variance: 
%
\begin{equation}
	 \sigma_{\min}^2 \leq \sigma_{i,j}^2 \leq \sigma^2, \qquad 1\leq i,j\leq n. 
	 \label{eq:sigma-min-definition}
\end{equation}
%
\begin{theorem}
\label{thm:matrix-distribution-complete}
%
Suppose that the assumptions of Theorem~\ref{thm:UsgnH-Ustar-MUstar-general} hold. For any $1\leq i,j\leq n$, set
%
\begin{align}
	v_{i,j}^{\star}=
		\begin{cases}
			\sum_{l=1}^{n}\sigma_{i,l}^{2}P^{\star 2}_{l,j}+\sum_{l=1}^{n}P^{\star 2}_{i,l}\sigma_{l,j}^{2}+2\sigma_{i,j}^{2}P^{\star}_{i,i}P^{\star}_{j,j}, & \text{if }i\neq j,\\
			4\sum_{l=1}^{n}\sigma_{i,l}^{2}P^{\star 2}_{l,i}, & \text{if }i=j. 
		\end{cases}
	\label{eq:variance-formula-Mij-lem}
\end{align}
%
Assume that $\sigma/\sigma_{\min}=O(1)$, and that
%
\begin{equation}
	\frac{\big\|\bm{U}_{j,\cdot}^{\star}\big\|_{2}^{2}+\big\|\bm{U}_{i,\cdot}^{\star}\big\|_{2}^{2}}{\big\|\bm{U}^{\star}\big\|_{\mathrm{F}}^{2}}
	\gtrsim \frac{B^{2}\kappa^{2}\mu^{2}r^{2}\log^{2}n}{\sigma^{2}n^{2}}+\frac{\sigma^{2}\mu^{2}r\kappa^{4}\log^{3}n}{(\lambda_{r}^{\star})^{2}}.
	%\frac{B^{2}\mu^{2}r\log n}{n^{2}\sigma^{2}}+\frac{\sigma^{2}\mu^{2}r\kappa^{4}\log^{3}n}{(\lambda_{r}^{\star})^{2}}+\frac{\kappa^{2}\mu^{2}r^{2}\log^{2}n}{n^{2}}.
	\label{eq:Ui-Uj-condition-thm-distribution-1}
\end{equation}
%
Then the estimator \eqref{eq:estimator-M-hat-distribution} obeys
%
\[
	\sup_{z\in\mathbb{R}}\left|\mathbb{P}\left(\widehat{M}_{i,j}-M_{i,j}^{\star}\leq z\sqrt{v_{i,j}^{\star}}\,\right)-\Phi(z)\right|=o(1) ,
\]
where $\Phi(\cdot)$ represents the cumulative density function (CDF) of the standard Gaussian distribution.
% 
\end{theorem}
%
The proof of this theorem is postponed to Section~\ref{sec:proof-thm:matrix-distribution-complete}. 
In a nutshell, Theorem~\ref{thm:matrix-distribution-complete} tells us that $\widehat{\bm{M}}$ is a nearly unbiased estimator of the truth $\bm{M}^{\star}$, as long as the signal strength---as captured by $\|\bm{U}^{\star}_{i,\cdot}\|_2$ and $\|\bm{U}^{\star}_{j,\cdot}\|_2$ when estimating the $(i,j)$-th entry---is sufficiently large (cf.~\eqref{eq:Ui-Uj-condition-thm-distribution-1}).  
The resulting estimation error in each entry is well approximated by a zero-mean Gaussian random variable, 
whose variance can be determined in a tractable fashion.  
As can be easily verified, the variance $v_{i,j}^{\star}$ is precisely the variance of the $(i,j)$-th entry of $\bm{E}\bm{U}^{\star}\bm{U}^{\star\top}+\bm{U}^{\star}\bm{U}^{\star\top}\bm{E}$ (as singled out in \eqref{eq:approx-ULambdaU-Mstar-discuss}). 
%In stark contrast to classical large-sample asymptotic analysis \citep{van2000asymptotic}, 
The above distributional theory is non-asymptotic, 
which lends itself well to high-dimensional applications.  





\subsection{Inference and uncertainty quantification}
\label{sec:UQ-general}



The Gaussian approximation unveiled in Theorem~\ref{thm:matrix-distribution-complete},
which is dictated by a single parameter ${v}_{i,j}^{\star}$, 
paves the way for statistical inference and uncertainty quantification tailored to this model. 
In order to construct a valid confidence interval for each entry of $\bm{M}^{\star}$, 
everything boils down to identifying an estimator that approximates the variance parameter ${v}_{i,j}^{\star}$, ideally in a data-driven yet faithful manner. 

In view of the variance characterization \eqref{eq:variance-formula-Mij-lem}, 
computing ${v}_{i,j}^{\star}$ requires information about both the noise variances $\{\sigma_{i,j}^2\}_{1\leq i,j\leq n}$ and the projection matrix $\bm{P}^{\star}$ (cf.~\eqref{eq:defn-P-UU-T}). However, estimating the noise variances is in general statistically infeasible, given that we only have access to a single observation (i.e., $M_{i,j}$) related to each individual variance $\sigma_{i,j}^2$.  Fortunately, the variance  ${v}_{i,j}^{\star}$ involves only the summation or equivalently the average of these individual variances, whose stochastic errors will be averaged out.  This leads us to the following surrogate
%
\begin{equation}
	\widetilde{v}_{i,j}=\begin{cases}
		\sum_{l=1}^{n}{E}_{i,l}^{2}{P}_{l,j}^{\star 2}+\sum_{l=1}^{n}{P}_{i,l}^{\star 2}{E}_{l,j}^{2}+2{E}_{i,j}^{2}{P}^{\star}_{i,i}{P}_{j,j}^{\star}, & \text{if }i\neq j,\\
	4\sum_{l=1}^{n}{E}_{i,l}^{2}{P}_{l,i}^{\star 2}, & \text{if }i=j, 
\end{cases}
	\label{eq:v-ij-plugin-surrogate-123}
\end{equation}
%
which is clearly an unbiased estimator of ${v}_{i,j}^{\star}$. 
Given the statistical independence of $\{E_{i,j}\}_{i\geq j}$, 
we can expect to have $\widetilde{v}_{i,j}\approx v_{i,j}^{\star}$, owing to the concentration of measure.  



%we propose to replace $\sigma_{i,j}^2$ by the corresponding entry $E_{i,j}^2$, or more realistically, an estimator $\widehat{E}_{i,j}^2$ of $E_{i,j}^2$. 



However, the above surrogate $\widetilde{v}_{i,j}$ remains practically incomputable, 
due to the absence of knowledge about both $\bm{E}$ and $\bm{P}^{\star}$.   
To address this issue, we propose the following plug-in estimator:
%
\begin{equation}
	\widehat{v}_{i,j}=\begin{cases}
	\sum_{l=1}^{n}\widehat{E}_{i,l}^{2}\widehat{P}_{l,j}^{2}+\sum_{l=1}^{n}\widehat{P}_{i,l}^{2}\widehat{E}_{l,j}^{2}+2\widehat{E}_{i,j}^{2}\widehat{P}_{i,i}\widehat{P}_{j,j}, & \text{if }i\neq j,\\
	4\sum_{l=1}^{n}\widehat{E}_{i,l}^{2}\widehat{P}_{l,i}^{2}, & \text{if }i=j, 
\end{cases}
	\label{eq:v-ij-plugin-estimator}
\end{equation}
%
where $\widehat{\bm{E}}=[\widehat{E}_{i,j}]_{1\leq i,j\leq n}$ and $\widehat{\bm{P}}=[\widehat{P}_{i,j}]_{1\leq i,j\leq n}$ stand for estimators of $\bm{E}$ and $\bm{P}^{\star}$, respectively. 
In particular, we employ the following specific estimators of $\bm{E}$ and $\bm{P}^{\star}$, again adopting the plug-in strategy:
%
\begin{subequations}
\label{eq:estimators-E-and-P}
\begin{align}
\widehat{\bm{E}} & \coloneqq\bm{M}-\bm{U}\bm{\Lambda}\bm{U}^{\top},\label{eq:estimator-noise-matrix-E}\\
\widehat{\bm{P}} & \coloneqq\bm{U}\bm{U}^{\top},\label{eq:estimator-projection-P}
\end{align}
\end{subequations}
%
where $\bm{U}$ and $\bm{\Lambda}$ are, as usual, computed via eigendecomposition of $\bm{M}$. 
For a prescribed coverage level $1-\alpha$ (with $0<\alpha <1$), we construct the following confidence interval for the $(i,j)$-th entry of $\bm{M}^{\star}$, 
motivated by the Gaussian approximation in Theorem~\ref{thm:matrix-distribution-complete}: 
%
\begin{align}
	\label{eq:confidence-interval-construction-general}
	\mathsf{CI}_{i,j}^{1-\alpha} \coloneqq \Big[\widehat{M}_{i,j}\pm\Phi^{-1}(1-\alpha/2)\sqrt{\widehat{v}_{i,j}}\,\Big].
\end{align}
%
Here and throughout, for any $b>0$, we let $[a\pm b]$ abbreviate the interval $[a-b,a+b]$, and we use $\Phi^{-1}(\cdot)$ to represent the inverse CDF of the standard Gaussian distribution. 


As encouraging news, the above construction of entrywise confidence intervals is provably valid with high probability, as revealed by the following theorem. The proof is postponed to Section~\ref{sec:proof-thm:CI-general}. 
%
\begin{theorem}
	\label{thm:CI-general}
	Consider the settings and assumptions in Section~\ref{sec:setup-general-theory},  
	%the assumptions of Theorem~\ref{thm:UsgnH-Ustar-MUstar-general} hold.
	and suppose that $\sigma / \sigma_{\min} = O(1)$,  $\kappa^{4}\mu^{2}r^{2}\log n\leq n$ and $\sigma \kappa\sqrt{n\log n}\lesssim |\lambda_{r}^{\star}|$. Consider any $1\leq i,j\leq n$, and assume that
	%
	\begin{align}
		\frac{\big\|\bm{U}_{j,\cdot}^{\star}\big\|_{2}^{2}+\big\|\bm{U}_{i,\cdot}^{\star}\big\|_{2}^{2}}{\big\|\bm{U}^{\star}\big\|_{\mathrm{F}}^{2}}\gtrsim\frac{B\kappa^{2}\mu^{2}r^{2}\log^{2}n}{\sigma n^{3/2}}+\frac{\sigma\mu^{2}r\kappa^{3}\log^{3}n}{|\lambda_{r}^{\star}|\sqrt{n}}.  
		\label{eq:Uj-Ui-lower-bound-thm}
	\end{align}
	%
	For any fixed coverage level $1-\alpha \in (0,1)$, 
	the confidence interval $\mathsf{CI}_{i,j}^{1-\alpha}$ constructed in \eqref{eq:confidence-interval-construction-general} obeys
	%
	\begin{align}
		%\left| 
		\mathbb{P} \Big( M_{i,j}^{\star} \in \mathsf{CI}_{i,j}^{1-\alpha}  
		\Big) 
		= 1-\alpha + o(1).
		% \right| =o(1).
	\end{align}
	%
\end{theorem}
%
%It is noteworthy that: 
Theorem~\ref{thm:CI-general} confirms that the confidence interval proposed above meets the prescribed coverage requirement,
provided that the associated signal strength is not too low (see \eqref{eq:Uj-Ui-lower-bound-thm}). 
In addition to its statistical validity, the proposed procedure enjoys several features that make it practically appealing: 
%
\begin{itemize}
	\item {\em Adaptive to unknown noise levels and distributions.} The above inference procedure is fully data-driven, which does not require prior knowledge about the noise levels or noise distributions. As alluded to previously, it is in general impossible to estimate the noise variance in each entry, and hence a data-driven yet valid approach is of critical value. 

	\item {\em Adaptive to heteroskedastic noise.} Our statistical guarantees hold without relying on homogeneity of noise components. In other words, this inference procedure automatically accommodates heteroskedastic noise, a scenario where the variance of the noise components might vary across different locations.   
\end{itemize}


Careful readers might remark that Theorem~\ref{thm:CI-general} is concerned with statistical inference for a single entry. 
Interestingly, the distributional theory presented in Section~\ref{sec:distribution-general} (see also Lemma~\ref{thm:matrix-distribution} in the proof of Theorem~\ref{thm:matrix-distribution-complete}) might also be instrumental in pursuing simultaneous inference, namely, the problem of constructing a confidence region that simultaneously accounts for more than one unknown entries. We omit such an extension for the sake of conciseness. 


%
%\end{proof}

%%
%\begin{align*}
% & \big\|\mathsf{sgn}(\bm{H})\bm{\Lambda}^{\star}\mathsf{sgn}(\bm{H})^{\top}-\bm{\Lambda}\big\|\\
% & \qquad\leq\big\|\mathsf{sgn}(\bm{H})\bm{\Lambda}^{\star}\mathsf{sgn}(\bm{H})^{\top}-\bm{H}\bm{\Lambda}^{\star}\bm{H}^{\top}\big\|+\big\|\bm{H}\bm{\Lambda}^{\star}\bm{H}^{\top}-\bm{\Lambda}\big\|\\
% & \qquad\lesssim\frac{\kappa\sigma^{2}n}{\lambda_{r}^{\star}}+\sigma\sqrt{r\log n}.
%\end{align*}



%
\section{Matrix completion}\label{sec:mc-l2}

A pressing challenge often encountered in data science applications is estimation and learning in the face of \emph{missing data}.
To elucidate how to tackle this challenge via spectral methods, we delve into the renowned matrix completion problem in this section,
followed by another application called tensor completion in Section~\ref{sec:tensor-completion}.


Imagine that one observes a small subset of the entries in a large unknown matrix
and seeks to fill in all missing entries. An archetypal example is collaborative filtering, where one aims to predict the users' preferences on a collection of products based on partially revealed user-product ratings.   See Figure~\ref{fig:completion} for an illustration.  The problem, often referred to as matrix completion, is apparently ill-posed in general, as there are (much) fewer measurements than the unknowns.

%This is indeed a factor model in Section~\ref{sec:PCA} with missing data.  

Fortunately, if the matrix of interest exhibits certain low-dimensional structure, then reliable recovery becomes feasible.
A commonly encountered example of this kind concerns the case when the target matrix enjoys a low-rank structure. Again, take collaborative filtering for example:  the user-product rating matrix  might be well explained by a relatively small number of latent factors connecting users' preferences with products' attributes, thus resulting in an approximately low-rank matrix.  Motivated by its fundamental importance,  recent years have witnessed a flurry of research activity in studying low-rank matrix completion \citep{candes2009exact, keshavan2010matrix, gross2011recovering};
see \citet{chen2018harnessing,davenport2016overview} for overviews of recent developments.
In the sequel, we present a simple yet effective approach enabled by the spectral method, originally proposed in~\citet{achlioptas2007fast,keshavan2010matrix}.

\begin{figure}
\begin{center}
\includegraphics[width=0.35\textwidth]{figures/matrix_completion}
\end{center}
\caption{Illustration of matrix completion, where each ``$\checkmark$'' stands for an observed entry and each ``$?$'' represents a missing entry.} \label{fig:completion}
\end{figure}


\subsection{Problem formulation and assumptions}
\label{sec:problem-formulation-MC}

Suppose that we are interested in estimating an $n_{1}\times n_{2}$  rank-$r$ matrix  $\bm{M}^{\star} = [{M}^{\star}_{i,j}]_{1\leq i\leq n_1, 1\leq j\leq n_2}$.
Without loss of generality, we assume
$$n_{1}\leq n_{2}.$$  Denote by  $\bm{M}^{\star}=\bm{U}^{\star}\bm{\Sigma}^{\star}\bm{V}^{\star\top}$ the SVD of $\bm{M}^{\star}$, where the columns of $\bm{U}^{\star}\in \mathbb{R}^{n_1\times r}$ (resp.~$\bm{V}^{\star}\in \mathbb{R}^{n_2\times r}$) are the left (resp.~right) singular vectors of $\bm{M}^{\star}$, and $\bm{\Sigma}^{\star}$ is a diagonal matrix whose diagonal entries are the singular values of $\bm{M}^{\star}$. We define the condition number of the matrix $\bm{M}^{\star}$ to be $\kappa \coloneqq \sigma_1(\bM^\star)/\sigma_r(\bM^\star)$.


To capture the presence of missing data, we introduce an index subset $\Omega\subseteq [n_1]\times [n_2]$,
such that each entry ${M}^{\star}_{i,j}$ is observed if and only if $(i,j)\in \Omega$.
The goal is to reconstruct the singular subspaces $\bm{U}^{\star}$ and $\bm{V}^{\star}$, as well as the  full matrix $\bm{M}^{\star}$, based on entries observed over the sampling set $\Omega$.



\paragraph{Random sampling.}
Apparently, not all sampling patterns admit reliable estimation. For instance, if $\Omega$ contains only entries in the top half of the matrix,
then there is in general no hope to predict the bottom half of the matrix. In order to allow for meaningful matrix completion,
this monograph focuses on a natural random observation model commonly adopted in the literature, as formulated below.
%
\begin{assumption}[Random sampling]
\label{assump:mc-bernoulli}
Each entry of $\bm{M}^{\star}$ is observed independently with probability $0<p<1$, namely,
each $(i,j)\in [ n_{1}] \times [n_{2}]$ is included in $\Omega$ independently with probability $p$.
\end{assumption}
%
Under this model, we shall view the expected number of observed entries---namely, $pn_1n_2$---as the sample size. In truth, as long as $p$ is not overly small,  the number of observed entries is expected to concentrate  around its mean $pn_1n_2$.



\paragraph{Incoherence conditions.}
Caution needs to be exercised, however,  that the random sampling model alone does not guarantee effective recovery of an arbitrary low-rank matrix $\bm{M}^{\star}$.
Consider, for example, the following rank-1 matrix $\bm{M}^{\star}$ containing a single nonzero entry:
%
\[
%{\small
\left[\begin{array}{cccc}
1 & 0 & \cdots & 0\\
0 & 0 & \cdots & 0\\
\vdots & \vdots & \ddots & \vdots\\
0 & 0 & \cdots & 0
\end{array}\right] .
%}
\]
%
If $p=o(1)$, then with probability $1-p=1-o(1)$, the sampling pattern will fail to include the nonzero entry $M_{1,1}^{\star}$,
thus ruling out the possibility of faithful matrix recovery.
Consequently,  one needs to make sure that the sampling pattern does not suppress too much useful information.
Towards this end, the pioneering work \citet{candes2009exact, candes2010NearOptimalMC} singled out an incoherence parameter that plays a vital role.
%
\begin{definition}
\label{assump:mc-incoherence}
The incoherence parameter $\mu$ of the matrix $\bm{M}^{\star}\in\mathbb{R}^{n_{1}\times n_{2}}$ is defined as
%
\begin{align*}
	\mu\coloneqq \max\left\{  \frac{ n_1 \|\bm{U}^{\star}\|_{2,\infty}^2 }{r} ,  \frac{ n_2 \|\bm{V}^{\star}\|_{2,\infty}^2 }{r} \right\} .
	%\\
	%\|\bm{V}^{\star}\|_{2,\infty} & \leq\sqrt{\mu/n_{2}}\,\|\bm{V}^{\star}\|_{\mathrm{F}}=\sqrt{\mu r/n_{2}}.
\end{align*}
%
\end{definition}
%
\begin{remark}
\label{remark:range-mu-MC}
%
Recognizing the following basic relation
%
\[
	\frac{r}{n_1}=\frac{1}{n_1} \|\bm{U}^{\star}\|_{\mathrm{F}}^2 \leq \|\bm{U}^{\star}\|_{2,\infty}^2 \leq \|\bm{U}^{\star}\|^2 = 1
\]
%
and an analogous one for $\bm{V}^{\star}$, we have $1\leq \mu \leq \max\{n_1,n_2\}/r= n_2/r$.
%
\end{remark}


In words, a small $\mu$ indicates that the energy of the singular vectors is spread out across different elements, namely,  the singular subspace of $\bm{M}^{\star}$ is not too ``aligned'' with any of the standard basis vectors,
thus ensuring that entrywise observations provide somewhat equalized information about the full spectrum of $\bm{M}^{\star}$.
The following lemma summarizes a few immediate consequences of this definition, with the proof deferred to Section~\ref{sec:proof-auxiliary-lemmas-MC-L2}.
%
\begin{lemma}\label{claim:mu-incoherence}
Assume that  $\bm{M}^{\star}\in\mathbb{R}^{n_{1}\times n_{2}}$ is $\mu$-incoherent. Then the following relations hold
%
\begin{subequations}
\label{subeq:mu-incoherence}
\begin{align}
	\|\bm{M}^{\star}\|_{2,\infty} & \leq\sqrt{\mu r/n_{1}}\, \big\| \bm{M}^{\star} \big\|; \quad
	\|\bm{M}^{\star\top}\|_{2,\infty}  \leq\sqrt{\mu r/n_{2}}\, \big\| \bm{M}^{\star} \big\|;\label{eq:mc-row-norm-upper}\\
	& \qquad\quad \|\bm{M}^{\star}\|_{\infty}  \leq\mu r  \big\| \bm{M}^{\star} \big\| /\sqrt{n_{1}n_{2}} . \label{eq:mc-entry-upper}
\end{align}
\end{subequations}
%
\end{lemma}



\paragraph{Additional notation.}

We find it convenient to introduce a Euclidean projection operator $\mathcal{P}_{\Omega}:\mathbb{R}^{n_{1}\times n_{2}}\mapsto\mathbb{R}^{n_{1}\times n_{2}}$ such that
%
\begin{equation}\label{defn:Pomega}
\big[\mathcal{P}_{\Omega}(\bm{A}) \big]_{i,j}=\begin{cases}
A_{i,j},\quad & \text{if }(i,j)\in\Omega \\
0, & \text{else}
\end{cases}
\end{equation}
%
for any matrix $\bm{A} =[A_{i,j}]\in\mathbb{R}^{n_{1}\times n_{2}}$.
With this notation in place, matrix completion amounts to recovering $\bm{M}^{\star}$ on the basis of $\mathcal{P}_{\Omega}(\bm{M}^{\star})$.




\subsection{Algorithm}\label{sec:mc-alg}

To apply the spectral method, the first step is to form a reasonable approximation $\bm{M}$ of the unknown matrix $\bm{M}^{\star}$.
By virtue of the random sampling model (cf.~Assumption~\ref{assump:mc-bernoulli}), a candidate approximation can be obtained from the observed data matrix via inverse probability weighting:
%
\begin{align}
	\bm{M}   \coloneqq p^{-1}\mathcal{P}_{\Omega}(\bm{M}^{\star}).
	\label{eq:rescaled-M-matrix-completion-l2}
\end{align}
%
The rationale is that $\bm{M}$ forms an unbiased estimate of the ground truth, namely, $$\mathbb{E}[\bm{M}]=\bm{M}^{\star},$$ where
the expectation is taken over the randomness in $\Omega$.


As a result, the proposed spectral method proceeds by computing the rank-$r$ SVD $\bm{U}\bm{\Sigma}\bm{V}^{\top}$ of the matrix $\bm{M}$ constructed in \eqref{eq:rescaled-M-matrix-completion-l2}, and employing $\bm{U}\in \mathbb{R}^{n_1\times r}$, $\bm{V}\in \mathbb{R}^{n_2\times r}$  and $\bm{U}\bm{\Sigma}\bm{V}^{\top}$ as estimates of $\bm{U}^{\star}$, $\bm{V}^{\star}$ and $\bm{M}^{\star}$, respectively.






\subsection{Performance guarantees}\label{sec:mc-l2-theory}

As before, whether the subspace $\bm{U}$ (resp.~$\bm{V}$) is close
to $\bm{U}^{\star}$ (resp.~$\bm{V}^{\star}$) relies crucially on
the size of the perturbation $\|\bm{M}-\bm{M}^{\star}\|$.
Therefore, we begin by developing an upper bound on this quantity; the proof is based on the matrix Bernstein inequality and is postponed to Section~\ref{sec:proof-auxiliary-lemmas-MC-L2}.

\begin{lemma}
\label{lemma:mc-noise-bound}
Consider the settings in Section~\ref{sec:problem-formulation-MC}. Suppose that $n_{2}p\geq C\mu r\log n_{2}$ for some
 constant $C>0$. Then with probability at least
	$1-O(n_{2}^{-10})$, the matrix $\bm{M}$ constructed in \eqref{eq:rescaled-M-matrix-completion-l2} obeys
%
\[
	\big\| \bm{M} -\bm{M}^{\star} \big\|
	\lesssim\sqrt{\frac{\mu r\log n_{2}}{n_{1}p}}\, \big\| \bm{M}^{\star} \big\|.
\]
%
\end{lemma}
%
With this perturbation bound in place, we are equipped to apply Wedin's $\sin\bm{\Theta}$ theorem to obtain the following results.  The condition on the sample size in Theorem~\ref{thm:mc-l2-subspace} is stronger than that in Lemma~\ref{lemma:mc-noise-bound}, as we need to control the eigengap in the following theorem.
%
\begin{theorem}
\label{thm:mc-l2-subspace}
Consider the settings in Section~\ref{sec:problem-formulation-MC}.
Suppose that $n_{1}p\geq C_{1}\kappa^{2}\mu r\log n_{2}$ for some sufficiently
large constant $C_{1}>0$. Then with probability exceeding $1-O(n_2^{-10})$,
%
\begin{align*}
\max\Big\{ \mathsf{dist}\left(\bm{U},\bm{U}^{\star}\right),\mathsf{dist}\left(\bm{V},\bm{V}^{\star}\right)\Big\}  & \lesssim\kappa\sqrt{\frac{\mu r \log n_{2}}{n_{1}p}} .
%;\\
%\max\Big\{ \mathsf{dist}_{\mathrm{F}}\left(\bm{U},\bm{U}^{\star}\right),\mathsf{dist}_{\mathrm{F}}\left(\bm{V},\bm{V}^{\star}\right)\Big\}  & \lesssim\kappa\sqrt{\frac{\mu r^{2} \log n_{2}}{n_{1}p}}.
\end{align*}
%
\end{theorem}
%
\begin{proof}
%	
As a direct consequence of Lemma \ref{lemma:mc-noise-bound}, one has
%
\[
	\big\| \bm{M} -\bm{M}^{\star} \big\|
	\lesssim\sqrt{\frac{\mu r\log n_{2}}{n_{1}p}} \,\big\| \bm{M}^{\star}\big\| \le \Big( 1 - \frac{1}{\sqrt{2}} \Big) \sigma_{r}(\bm{M}^{\star}),
\]
%
provided that $n_{1}p\geq C_{1}\kappa^{2}\mu r\log n_{2}$ for some  large enough constant $C_{1}>0$.
Apply Wedin's theorem (cf.~\eqref{eq:wedin-simpler-friendly}) and Lemma \ref{lemma:mc-noise-bound} to obtain
%
\begin{align*}
 & \max\Big\{ \mathsf{dist}\left(\bm{U},\bm{U}^{\star}\right),\mathsf{dist}\left(\bm{V},\bm{V}^{\star}\right) \Big\}
	\leq  \frac{2\big\|\bm{M} -\bm{M}^{\star}\big\|}{\sigma_{r}(\bm{M}^{\star})}
	%\\ % & \qquad\qquad
	%\lesssim\frac{\sqrt{\frac{\mu r\log n_{2}}{n_{1}p}}\,\big\|\bm{M}^{\star}\big\|}{\sigma_{r}(\bm{M}^{\star})}\asymp
	\lesssim \kappa\sqrt{\frac{\mu r\log n_{2}}{n_{1}p}}
\end{align*}
%
as claimed.
%The proof for the $\mathsf{dist}_{\mathrm{F}}(\cdot,\cdot)$ bounds follows from the same argument and is hence omitted.
%
\end{proof}


As an important implication of Theorem~\ref{thm:mc-l2-subspace}, once the sample size exceeds
$$pn_1n_2 \gg \kappa^{2}\mu r n_2  \log n_{2}, $$
then the spectral estimate achieves consistent estimation in the sense that
%
\[
	\max\Big\{ \mathsf{dist}\left(\bm{U},\bm{U}^{\star}\right),\mathsf{dist}\left(\bm{V},\bm{V}^{\star}\right)\Big\} = o(1).
\]
%
Given that $pn_1n_2\gtrsim \mu n_2 r \log n_2$ is an information-theoretic sampling requirement for reliable matrix completion when $r=o(n_1/\log n_2)$ \citep{candes2010NearOptimalMC},
Theorem~\ref{thm:mc-l2-subspace} confirms the near optimality of spectral methods---in terms of the scaling with $n_1$, $n_2$ and $p$---when it comes to consistent subspace estimation.



Before moving forward to the proof, we further characterize the statistical accuracy of $\bm{U}\bm{\Sigma}\bm{V}^{\top}$ in estimating the unknown matrix $\bm{M}^{\star}$. Accomplishing this only requires Lemma~\ref{lemma:mc-noise-bound}, without any need of the singular subspace perturbation theory. This result will also come in handy when we turn to discussing entrywise estimation accuracy in Chapter~\ref{cha:Linf-theory}.
%
\begin{theorem}
\label{thm:mc-l2-matrix}
Consider the settings in Section~\ref{sec:problem-formulation-MC}.
Suppose that $n_{2}p\geq C\mu r\log n_{2}$ for some sufficiently large constant
$C>0$. Then with probability at least $1-O(n_{2}^{-10})$, one has
%
\begin{align*}
%\|\bm{U}\bm{\Sigma}\bm{V}^{\top}-\bm{M}^{\star}\| & \lesssim\sqrt{\frac{\mu r}{n_{1}p}\log n_{2}}\,\sigma_{1}(\bm{M}^{\star});\\
\|\bm{U}\bm{\Sigma}\bm{V}^{\top}-\bm{M}^{\star}\|_{\mathrm{F}} & \lesssim\sqrt{\frac{\mu r^{2}\log n_{2}}{n_{1}p}}\, \big\| \bm{M}^{\star} \big\|.
\end{align*}
%
\end{theorem}
%
\begin{proof}
First, note
%
\begin{align*}
\|\bm{U}\bm{\Sigma}\bm{V}^{\top}-\bm{M}^{\star}\| &\leq\|\bm{U}\bm{\Sigma}\bm{V}^{\top}- \bm{M} \|+\|\bm{M} -\bm{M}^{\star}\| \leq2\|\bm{M}-\bm{M}^{\star}\|,
\end{align*}
%
where the first inequality comes from the triangle inequality, and the second inequality follows from the fact that $\bm{U}\bm{\Sigma}\bm{V}^{\top}$ is the best rank-$r$ approximation to $\bm{M}$,
i.e.,
$$\|\bm{U}\bm{\Sigma}\bm{V}^{\top}-\bm{M}\|=\min_{\bm{Z}: \mathsf{rank}(\bm{Z})\leq r}\|\bm{Z}-\bm{M}\| \leq \|\bm{M}-\bm{M}^{\star}\|. $$
Additionally, it is observed that $\bm{U}\bm{\Sigma}\bm{V}^{\top}-\bm{M}^{\star}$ has rank at most $2r$, which implies
%
\[
	\|\bm{U}\bm{\Sigma}\bm{V}^{\top}-\bm{M}^{\star}\|_{\mathrm{F}}\leq\sqrt{2r}\|\bm{U}\bm{\Sigma}\bm{V}^{\top}-\bm{M}^{\star}\| \leq 2\sqrt{2r} \big\|\bm{M} - \bm{M}^{\star} \big\|.
\]
%
This combined with Lemma~\ref{lemma:mc-noise-bound} immediately concludes the proof. \end{proof}

 


\subsection{Proof of auxiliary lemmas}
\label{sec:proof-auxiliary-lemmas-MC-L2}

\paragraph{Proof of Lemma~\ref{claim:mu-incoherence}.}
First of all, the  $\|\cdot\|_{2,\infty}$ norm of $\bm{M}^{\star}$
can be upper bounded by
%
\begin{equation*}
\|\bm{M}^{\star}\|_{2,\infty}=\|\bm{U}^{\star}\bm{\Sigma}^{\star}\bm{V}^{\star\top}\|_{2,\infty}\leq\|\bm{U}^{\star}\|_{2,\infty}\|\bm{\Sigma}^{\star}\| \, \|\bm{V}^{\star}\|\leq\sqrt{\frac{\mu r}{n_{1}}} \big\| \bm{M}^{\star} \big\|.
\end{equation*}
%
Here, the first inequality arises from the elementary bounds
$\|\bm{A}\bm{B}\|_{2,\infty}\leq\|\bm{A}\|_{2,\infty}\|\bm{B}\|$
and $\|\bm{A}\bm{B}\|\leq\|\bm{A}\|\, \|\bm{B}\|$, whereas the last relation
uses Definition~\ref{assump:mc-incoherence}, the orthonormality
of $\bm{V}^{\star}$, and identifies $\|\bm{\Sigma}^{\star}\|$ with
$\| \bm{M}^{\star} \|$. The  bound on $\|\bm{M}^{\star\top}\|_{2,\infty}$ can be derived analogously and is omitted for brevity.



In addition, the matrix $\bm{M}^{\star}$ is elementwise bounded by
%
\begin{equation*}
\|\bm{M}^{\star}\|_{\infty}=\|\bm{U}^{\star}\bm{\Sigma}^{\star}\bm{V}^{\star\top}\|_{\infty}\leq\|\bm{U}^{\star}\|_{2,\infty}\|\bm{V}^{\star}\|_{2,\infty} \|\bm{\Sigma}^{\star}\|
	\leq\frac{\mu r}{\sqrt{n_{1}n_{2}}} \big\| \bm{M}^{\star} \big\|.
\end{equation*}
%
Here, the first inequality follows from the fact $\|\bm{A}\bm{B}^{\top}\|_{\infty}\leq\|\bm{A}\|_{2,\infty}\|\bm{B}\|_{2,\infty}$
and the aforementioned one $\|\bm{A}\bm{B}\|_{2,\infty}\leq\|\bm{A}\|_{2,\infty}\|\bm{B}\|$,
while the last inequality again relies on Definition~\ref{assump:mc-incoherence}.
%\end{proof}


\paragraph{Proof of Lemma~\ref{lemma:mc-noise-bound}.} Note that the matrix $\bm{E}\coloneqq p^{-1}\mathcal{P}_{\Omega}(\bm{M}^{\star})-\bm{M}^{\star}$
can be expressed as the sum of $n_{1}n_{2}$ i.i.d.~random matrices
%
\[
  \frac{1}{p}  \mathcal{P}_{\Omega}(\bm{M}^{\star})-\bm{M}^{\star}=\sum_{i=1}^{n_{1}}\sum_{j=1}^{n_{2}}\underbrace{\big(p^{-1} \delta_{i,j} -1 \big)M_{i,j}^{\star}\bm{e}_{i}\bm{e}_{j}^{\top}}_{\eqqcolon\bm{X}_{i,j}}.
\]
%
Here, $\delta_{i,j}$ (which indicates whether the $(i,j)$-th entry is observed) follows an independent Bernoulli distribution
with parameter~$p$, and $\bm{e}_{i}$ stands for the $i$-th standard
basis vector of appropriate dimensions. It is easily seen that for each $(i,j)$,
%
\[
	\mathbb{E}[\bm{X}_{i,j}]=\bm{0}\quad\text{and}\quad\|\bm{X}_{i,j}\|\leq  \frac{1}{p} \|\bm{M}^{\star}\|_{\infty}\leq\frac{\mu r}{p\sqrt{n_{1}n_{2}}} \big\| \bm{M}^{\star} \big\|,
\]
%
where the last relation results from the entrywise upper bound (\ref{eq:mc-entry-upper})
on $\bm{M}^{\star}$. In order to apply the matrix Bernstein inequality (cf.~Corollary~\ref{thm:matrix-Bernstein-friendly}),
we need to control the variance statistic
%
\[
	v\coloneqq\max \Big\{ \Big\Vert \sum\nolimits _{i,j}\mathbb{E}\left[\bm{X}_{i,j}\bm{X}_{i,j}^{\top}\right]\Big\Vert ,\Big\Vert \sum\nolimits _{i,j}\mathbb{E}\left[\bm{X}_{i,j}^{\top}\bm{X}_{i,j}\right] \Big\Vert \Big\} .
\]
%
Regarding the first variance term, we have
%
\begin{align*}
\sum\nolimits _{i,j}\mathbb{E}\left[\bm{X}_{i,j}\bm{X}_{i,j}^{\top}\right]
	& =\sum\nolimits _{i,j}\mathbb{E}\left[ \big( p^{-1} \delta_{i,j} -1 \big)^{2}(M_{i,j}^{\star})^{2}\bm{e}_{i}\bm{e}_{j}^{\top}\bm{e}_{j}\bm{e}_{i}^{\top}\right]\\
 & =\frac{1-p}{p}\sum\nolimits _{i,j}\big(M_{i,j}^{\star}\big)^{2}\bm{e}_{i}\bm{e}_{i}^{\top}=\frac{1-p}{p}\sum_{i=1}^{n_{1}}\|\bm{M}_{i,\cdot}^{\star}\|_{2}^{2}\bm{e}_{i}\bm{e}_{i}^{\top} \\
	& \preceq \frac{1-p}{p}  \|\bm{M}^{\star}\|_{2,\infty}^{2} \bm{I}_{n_1} \preceq \frac{\mu r}{n_1p}  \|\bm{M}^{\star}\|^{2} \,\bm{I}_{n_1} .
\end{align*}
%
Here, the first identity arises from the definition of $\bm{X}_{i,j}$,
the second one calculates the variance of Bernoulli random variables,
and the last line relies on the upper bound~(\ref{eq:mc-row-norm-upper}).
Similarly, the second term in the variance statistic enjoys the following characterization:
%
\[
	\sum\nolimits _{i,j}\mathbb{E}\left[\bm{X}_{i,j}^{\top}\bm{X}_{i,j}\right] \preceq \frac{\mu r}{n_2p} \|\bm{M}^{\star}\|^{2} \,\bm{I}_{n_2}.
\]
%
Taking the above  relations together and recalling that $n_1\leq n_2$ give
%
\begin{align*}
v & %\leq\frac{1}{p}\max\left\{ \|\bm{M}^{\star}\|_{2,\infty}^{2},\|\bm{M}^{\star\top}\|_{2,\infty}^{2}\right\}
	% \leq\frac{1}{p}\left\{ \frac{\mu r}{n_{1}} \big\| \bm{M}^{\star} \big\|^2,  \frac{\mu r}{n_{2}} \big\| \bm{M}^{\star} \big\|^2 \right\} \\
	\leq\frac{\mu r}{n_{1}p} \big\| \bm{M}^{\star} \big\|^2.
\end{align*}
% 


With the above bounds in place, invoking matrix Bernstein (see Corollary~\ref{thm:matrix-Bernstein-friendly}) reveals that: with probability at least $1-O(n_2^{-10})$,
%
\begin{align*}
	\| \bm{E}\| & \lesssim \sqrt{\frac{\mu r \| \bm{M}^{\star} \|^2 \log n_{2}}{n_{1}p} }+\frac{\mu r \| \bm{M}^{\star} \| \log n_{2}}{p\sqrt{n_{1}n_{2}}}
	\asymp \sqrt{\frac{\mu r \| \bm{M}^{\star} \|^2 \log n_{2}}{n_{1}p} } ,
\end{align*}
%
where the last inequality is valid as long as  $n_{2}p \gtrsim \mu r\log n_{2}$.







\section{Application: Confidence intervals for matrix completion}
\label{sec:CI-matrix-completion}

As an illustration of the applicability of the inference procedure described in Section~\ref{sec:UQ-general}, 
we develop concrete consequences of Theorem~\ref{thm:CI-general} in application to noisy matrix completion---an extension of the formulation
in Section~\ref{sec:mc-l2} to noisy settings. 


\paragraph{Model: noisy matrix completion.} 

Suppose that we are asked to reconstruct a symmetric rank-$r$ matrix
 $\bm{M}^{\star}=[M_{i,j}^{\star}]_{1\leq k,l \leq n}\in \mathbb{R}^{n\times n}$  
 with eigendecomposition $\bm{M}^{\star}=\bm{U}^{\star}\bm{\Lambda}^{\star}\bm{U}^{\star\top}$.
 We only get to acquire noisy observations of a subset of the entries of $\bm{M}^{\star}$; more precisely, 
there exists a sampling set $\Omega \subseteq [n]\times [n]$ such that we observe
%
\begin{align}
	M_{k,l}^{\star} + \eta_{k,l} \qquad & \text{if } (k,l)\in \Omega. 
	% M_{i,j} = \begin{cases} , \qquad & \text{if } (i,j)\in \Omega, \\ 0, &\text{else.}  \end{cases}
\end{align}
%
Here, $\{\eta_{k,l}\mid k\geq l\}$ denotes independent Gaussian noise obeying  
%
\begin{align}
	\eta_{k,l} = \eta_{l,k} \overset{\mathrm{i.i.d.}}{\sim} \mathcal{N}(0, \sigma_{\eta}^2), \qquad k\geq l. 
\end{align}
%
As before, we focus on the random sampling model such that each location $(k,l)$ with $k\geq l$ is included in the sampling set $\Omega$ independently with probability $p$. Further,
 assume that $\bm{M}^{\star}$ has eigenvalues obeying \eqref{eq:assumption-lambda-r-positive-setup}, condition number $\kappa$ (cf.~\eqref{eq:defn-kappa}), and  incoherence parameter $\mu$ (cf.~\eqref{eq:defn-Ustar-incoherence}). 
Can we build a confidence interval for each entry $M_{i,j}^{\star}$, on the basis of the output of the spectral method?  






\paragraph{Computing entrywise confidence intervals.} 
%
In order to apply the inference procedure in Section~\ref{sec:UQ-general}, it suffices to determine the data matrix $\bm{M}=[M_{i,j}]_{1\leq i,j\leq n}$, which can be selected as usual. Specifically, a possible inference procedure proceeds as follows:
%
\begin{itemize}
	\item  Set $\bm{M}$ such that for any $1\leq i,j \leq n$, 
%
\begin{align}
	M_{i,j} = \begin{cases} \frac{1}{p} \big( M_{i,j}^{\star} + \eta_{i,j} \big), \qquad & \text{if } (i,j)\in \Omega, \\ 0, &\text{else}, \end{cases}
\end{align}
%
which clearly obeys $\mathbb{E}[\bm{M}]=\bm{M}^{\star}$. 

	\item Compute the estimate $\widehat{\bm{M}}$ (cf.~\eqref{eq:estimator-M-hat-distribution}) via the spectral method.

	\item For a given coverage level $1-\alpha$ and a given pair $(i,j)$, 
		construct the confidence interval $\mathsf{CI}_{i,j}^{1-\alpha}$ according to \eqref{eq:confidence-interval-construction-general}, 
		with auxiliary parameters provided in \eqref{eq:v-ij-plugin-estimator} and \eqref{eq:estimators-E-and-P}. 

\end{itemize}
 


\paragraph{Performance guarantees and implications.}
When specialized to noisy matrix completion, our inference theory in Theorem~\ref{thm:CI-general} leads to the following statistical guarantees. 
%
\begin{theorem}
	\label{thm:CI-noisy-matrix-completion}
	Consider the noisy matrix completion setting in this section. 
	Suppose that $\kappa^{4}\mu^{2}r^{2}\log n\leq n$, 
	%$\sigma \kappa\sqrt{n\log n}\lesssim |\lambda_{r}^{\star}|$ and 
	$\frac{\max_{k,l}|M_{k,l}^{\star}|}{\min_{k,l}|M_{k,l}^{\star}|}=O(1)$, 
	%
	\begin{equation}
		np\gtrsim \kappa^{4}r\log n\quad\text{and}\quad\sigma_{\eta}\kappa\sqrt{\frac{n\log n}{p}}\lesssim|\lambda_{r}^{\star}|.  
		\label{eq:Mstar-noise-requirement-noisy-MC}
	\end{equation}
	%	
	Consider any $1\leq i,j\leq n$, and assume that
	%
\begin{align}
	\frac{\big\|\bm{U}_{j,\cdot}^{\star}\big\|_{2}^{2}+\big\|\bm{U}_{i,\cdot}^{\star}\big\|_{2}^{2}}{\big\|\bm{U}^{\star}\big\|_{\mathrm{F}}^{2}} & \gtrsim\frac{\mu^{2}r^{2}\kappa^{4}\log^{3}n}{n\sqrt{np}}+\frac{\sigma_{\eta}\mu^{2}r\kappa^{3}\log^{3}n}{|\lambda_{r}^{\star}|\sqrt{np}}.
	\label{eq:Ui-Uj-lower-bound-thm-noisy-MC-146}
\end{align}
%
	For any fixed coverage level $1-\alpha \in (0,1)$, 
	the confidence interval $\mathsf{CI}_{i,j}^{1-\alpha}$ constructed in \eqref{eq:confidence-interval-construction-general} obeys
	%
	\begin{align}
		\mathbb{P} \Big( M_{i,j}^{\star} \in \mathsf{CI}_{i,j}^{1-\alpha}  
		\Big) 
		= 1-\alpha + o(1).
	\end{align}
	%
\end{theorem}
%


In order to help interpret the applicable range of Theorem~\ref{thm:CI-noisy-matrix-completion}, 
let us focus on the simple scenario with $\kappa, \mu, r\asymp 1$ to simplify discussion. 

\begin{itemize}

	\item First of all, Condition~\eqref{eq:Mstar-noise-requirement-noisy-MC} can be simplified as
	%
	\[
		np  \gtrsim  \log n\quad\text{and}\quad \sigma_{\eta}\sqrt{\frac{n\log n}{p}}\lesssim |\lambda_r^{\star}|. 
	\]
	%
	The first condition on the sampling size coincides with the fundamental requirement even if the goal is merely to enable reliable estimation \citep{candes2010NearOptimalMC}, 
	whereas the second condition on the signal-to-noise ratio is also necessary---up to some log factor---to ensure an estimation quality better than that of a random guess \citep[Theorem~3.3]{cai2019subspace}.  


	\item Next, we move on to interpret the other condition \eqref{eq:Ui-Uj-lower-bound-thm-noisy-MC-146} imposed in our theory, which simplifies to
	%
	\[
		\frac{\big\|\bm{U}_{j,\cdot}^{\star}\big\|_{2}^{2}+\big\|\bm{U}_{i,\cdot}^{\star}\big\|_{2}^{2}}{\big\|\bm{U}^{\star}\big\|_{\mathrm{F}}^{2}}
		\gtrsim\frac{\log^{3}n}{n\sqrt{np}}+\frac{\sigma_{\eta}\log^{3}n}{|\lambda_{r}^{\star}|\sqrt{np}}.
	\]
	%
	Let us consider the most challenging case where $np \gtrsim  \mathrm{poly}\log n$ and $\sigma_{\eta}\sqrt{\frac{n\mathrm{poly}\log n}{p}}\lesssim |\lambda_r^{\star}|$ (for some sufficiently large poly-log factor). 
	In such a case, the above condition only requires 
	%
	\[
		\frac{\big\|\bm{U}_{j,\cdot}^{\star}\big\|_{2}^{2}+\big\|\bm{U}_{i,\cdot}^{\star}\big\|_{2}^{2}}{\big\|\bm{U}^{\star}\big\|_{\mathrm{F}}^{2}}\gtrsim\frac{1}{n\mathrm{poly}\log n},
	\]
	%
	indicating that the associated signal power $\|\bm{U}_{j,\cdot}^{\star}\|_{2}^{2}+\|\bm{U}_{i,\cdot}^{\star}\|_{2}^{2}$ 
		is allowed to be much smaller than the average signal power across all rows (which can be captured by 
		$\|\bm{U}^{\star}\|_{\mathrm{F}}^{2}/n$). 

\end{itemize}
%
In a nutshell, the validity of our inference procedure is ensured for broad settings. 
Additionally, we have conducted a series of numerical experiments to examine the entrywise distributions of $\bm{M}$. As illustrated in Figure~\ref{fig:MC-inference-numerics}, 
the normalized estimation error $(\widehat{v}_{i,j})^{-1/2} (\widehat{M}_{i,j} - M_{i,j}^{\star})$ is close in distribution to a standard Gaussian random variable, 
which corroborates our theory on the confidence interval construction. 



Before concluding, we would like to remark that: 
while the distributional theory for spectral methods allows for valid construction of confidence intervals for an unseen entry,  
it is oftentimes not among the most effective statistical inference procedures that one can put forward. 
There exist other alternatives that are provably more efficient, including but not limited to inference procedures based on convex relaxation and nonconvex optimization \citep{chen2019inference,xia2021statistical}, and the ones based on more refined spectral methods \citep{yan2021inference,chernozhukov2021inference}. 


\begin{figure}

	\begin{tabular}{cc}
		\includegraphics[width=0.45\textwidth]{figures/mc_inference_histogram_plot.pdf} & \includegraphics[width=0.45\textwidth]{figures/mc_inference_QQ_plot.pdf} \tabularnewline
		(a)  & (b) \tabularnewline
	\end{tabular}
	
	\caption{Entrywise numerical distribution for noisy matrix completion. We generate $\bm{M}^{\star}=\bm{U}^{\star}\bm{U}^{\star\top}$ with $\bm{U}^{\star}$ being a random orthonormal matrix, and analyze the behavior of $Z_{1,2}= (\widehat{v}_{1,2})^{-0.5} ( \widehat{M}_{1,2}-M_{1,2}^{\star} )$, 
		where $\widehat{\bm{M}}$ is defined in \eqref{eq:estimator-M-hat-distribution} and $\widehat{v}_{i,j}$ is defined in \eqref{eq:v-ij-plugin-estimator}.
		The results are reported for 500 Monte Carlo trials when $p=0.3$, $n=1000$, $\sigma_{\eta} = 10^{-4}$, and $r=3$.  
		 (a) Histogram of the empirical distribution of $Z_{1,2}$; (b)  Q-Q (quantile-quantile) plot of $Z_{1,2}$ vs.~the standard normal distribution. \label{fig:MC-inference-numerics}}  
\end{figure}



%
\section{Matrix completion}\label{sec:mc-l2}

A pressing challenge often encountered in data science applications is estimation and learning in the face of \emph{missing data}.
To elucidate how to tackle this challenge via spectral methods, we delve into the renowned matrix completion problem in this section,
followed by another application called tensor completion in Section~\ref{sec:tensor-completion}.


Imagine that one observes a small subset of the entries in a large unknown matrix
and seeks to fill in all missing entries. An archetypal example is collaborative filtering, where one aims to predict the users' preferences on a collection of products based on partially revealed user-product ratings.   See Figure~\ref{fig:completion} for an illustration.  The problem, often referred to as matrix completion, is apparently ill-posed in general, as there are (much) fewer measurements than the unknowns.

%This is indeed a factor model in Section~\ref{sec:PCA} with missing data.  

Fortunately, if the matrix of interest exhibits certain low-dimensional structure, then reliable recovery becomes feasible.
A commonly encountered example of this kind concerns the case when the target matrix enjoys a low-rank structure. Again, take collaborative filtering for example:  the user-product rating matrix  might be well explained by a relatively small number of latent factors connecting users' preferences with products' attributes, thus resulting in an approximately low-rank matrix.  Motivated by its fundamental importance,  recent years have witnessed a flurry of research activity in studying low-rank matrix completion \citep{candes2009exact, keshavan2010matrix, gross2011recovering};
see \citet{chen2018harnessing,davenport2016overview} for overviews of recent developments.
In the sequel, we present a simple yet effective approach enabled by the spectral method, originally proposed in~\citet{achlioptas2007fast,keshavan2010matrix}.

\begin{figure}
\begin{center}
\includegraphics[width=0.35\textwidth]{figures/matrix_completion}
\end{center}
\caption{Illustration of matrix completion, where each ``$\checkmark$'' stands for an observed entry and each ``$?$'' represents a missing entry.} \label{fig:completion}
\end{figure}


\subsection{Problem formulation and assumptions}
\label{sec:problem-formulation-MC}

Suppose that we are interested in estimating an $n_{1}\times n_{2}$  rank-$r$ matrix  $\bm{M}^{\star} = [{M}^{\star}_{i,j}]_{1\leq i\leq n_1, 1\leq j\leq n_2}$.
Without loss of generality, we assume
$$n_{1}\leq n_{2}.$$  Denote by  $\bm{M}^{\star}=\bm{U}^{\star}\bm{\Sigma}^{\star}\bm{V}^{\star\top}$ the SVD of $\bm{M}^{\star}$, where the columns of $\bm{U}^{\star}\in \mathbb{R}^{n_1\times r}$ (resp.~$\bm{V}^{\star}\in \mathbb{R}^{n_2\times r}$) are the left (resp.~right) singular vectors of $\bm{M}^{\star}$, and $\bm{\Sigma}^{\star}$ is a diagonal matrix whose diagonal entries are the singular values of $\bm{M}^{\star}$. We define the condition number of the matrix $\bm{M}^{\star}$ to be $\kappa \coloneqq \sigma_1(\bM^\star)/\sigma_r(\bM^\star)$.


To capture the presence of missing data, we introduce an index subset $\Omega\subseteq [n_1]\times [n_2]$,
such that each entry ${M}^{\star}_{i,j}$ is observed if and only if $(i,j)\in \Omega$.
The goal is to reconstruct the singular subspaces $\bm{U}^{\star}$ and $\bm{V}^{\star}$, as well as the  full matrix $\bm{M}^{\star}$, based on entries observed over the sampling set $\Omega$.



\paragraph{Random sampling.}
Apparently, not all sampling patterns admit reliable estimation. For instance, if $\Omega$ contains only entries in the top half of the matrix,
then there is in general no hope to predict the bottom half of the matrix. In order to allow for meaningful matrix completion,
this monograph focuses on a natural random observation model commonly adopted in the literature, as formulated below.
%
\begin{assumption}[Random sampling]
\label{assump:mc-bernoulli}
Each entry of $\bm{M}^{\star}$ is observed independently with probability $0<p<1$, namely,
each $(i,j)\in [ n_{1}] \times [n_{2}]$ is included in $\Omega$ independently with probability $p$.
\end{assumption}
%
Under this model, we shall view the expected number of observed entries---namely, $pn_1n_2$---as the sample size. In truth, as long as $p$ is not overly small,  the number of observed entries is expected to concentrate  around its mean $pn_1n_2$.



\paragraph{Incoherence conditions.}
Caution needs to be exercised, however,  that the random sampling model alone does not guarantee effective recovery of an arbitrary low-rank matrix $\bm{M}^{\star}$.
Consider, for example, the following rank-1 matrix $\bm{M}^{\star}$ containing a single nonzero entry:
%
\[
%{\small
\left[\begin{array}{cccc}
1 & 0 & \cdots & 0\\
0 & 0 & \cdots & 0\\
\vdots & \vdots & \ddots & \vdots\\
0 & 0 & \cdots & 0
\end{array}\right] .
%}
\]
%
If $p=o(1)$, then with probability $1-p=1-o(1)$, the sampling pattern will fail to include the nonzero entry $M_{1,1}^{\star}$,
thus ruling out the possibility of faithful matrix recovery.
Consequently,  one needs to make sure that the sampling pattern does not suppress too much useful information.
Towards this end, the pioneering work \citet{candes2009exact, candes2010NearOptimalMC} singled out an incoherence parameter that plays a vital role.
%
\begin{definition}
\label{assump:mc-incoherence}
The incoherence parameter $\mu$ of the matrix $\bm{M}^{\star}\in\mathbb{R}^{n_{1}\times n_{2}}$ is defined as
%
\begin{align*}
	\mu\coloneqq \max\left\{  \frac{ n_1 \|\bm{U}^{\star}\|_{2,\infty}^2 }{r} ,  \frac{ n_2 \|\bm{V}^{\star}\|_{2,\infty}^2 }{r} \right\} .
	%\\
	%\|\bm{V}^{\star}\|_{2,\infty} & \leq\sqrt{\mu/n_{2}}\,\|\bm{V}^{\star}\|_{\mathrm{F}}=\sqrt{\mu r/n_{2}}.
\end{align*}
%
\end{definition}
%
\begin{remark}
\label{remark:range-mu-MC}
%
Recognizing the following basic relation
%
\[
	\frac{r}{n_1}=\frac{1}{n_1} \|\bm{U}^{\star}\|_{\mathrm{F}}^2 \leq \|\bm{U}^{\star}\|_{2,\infty}^2 \leq \|\bm{U}^{\star}\|^2 = 1
\]
%
and an analogous one for $\bm{V}^{\star}$, we have $1\leq \mu \leq \max\{n_1,n_2\}/r= n_2/r$.
%
\end{remark}


In words, a small $\mu$ indicates that the energy of the singular vectors is spread out across different elements, namely,  the singular subspace of $\bm{M}^{\star}$ is not too ``aligned'' with any of the standard basis vectors,
thus ensuring that entrywise observations provide somewhat equalized information about the full spectrum of $\bm{M}^{\star}$.
The following lemma summarizes a few immediate consequences of this definition, with the proof deferred to Section~\ref{sec:proof-auxiliary-lemmas-MC-L2}.
%
\begin{lemma}\label{claim:mu-incoherence}
Assume that  $\bm{M}^{\star}\in\mathbb{R}^{n_{1}\times n_{2}}$ is $\mu$-incoherent. Then the following relations hold
%
\begin{subequations}
\label{subeq:mu-incoherence}
\begin{align}
	\|\bm{M}^{\star}\|_{2,\infty} & \leq\sqrt{\mu r/n_{1}}\, \big\| \bm{M}^{\star} \big\|; \quad
	\|\bm{M}^{\star\top}\|_{2,\infty}  \leq\sqrt{\mu r/n_{2}}\, \big\| \bm{M}^{\star} \big\|;\label{eq:mc-row-norm-upper}\\
	& \qquad\quad \|\bm{M}^{\star}\|_{\infty}  \leq\mu r  \big\| \bm{M}^{\star} \big\| /\sqrt{n_{1}n_{2}} . \label{eq:mc-entry-upper}
\end{align}
\end{subequations}
%
\end{lemma}



\paragraph{Additional notation.}

We find it convenient to introduce a Euclidean projection operator $\mathcal{P}_{\Omega}:\mathbb{R}^{n_{1}\times n_{2}}\mapsto\mathbb{R}^{n_{1}\times n_{2}}$ such that
%
\begin{equation}\label{defn:Pomega}
\big[\mathcal{P}_{\Omega}(\bm{A}) \big]_{i,j}=\begin{cases}
A_{i,j},\quad & \text{if }(i,j)\in\Omega \\
0, & \text{else}
\end{cases}
\end{equation}
%
for any matrix $\bm{A} =[A_{i,j}]\in\mathbb{R}^{n_{1}\times n_{2}}$.
With this notation in place, matrix completion amounts to recovering $\bm{M}^{\star}$ on the basis of $\mathcal{P}_{\Omega}(\bm{M}^{\star})$.




\subsection{Algorithm}\label{sec:mc-alg}

To apply the spectral method, the first step is to form a reasonable approximation $\bm{M}$ of the unknown matrix $\bm{M}^{\star}$.
By virtue of the random sampling model (cf.~Assumption~\ref{assump:mc-bernoulli}), a candidate approximation can be obtained from the observed data matrix via inverse probability weighting:
%
\begin{align}
	\bm{M}   \coloneqq p^{-1}\mathcal{P}_{\Omega}(\bm{M}^{\star}).
	\label{eq:rescaled-M-matrix-completion-l2}
\end{align}
%
The rationale is that $\bm{M}$ forms an unbiased estimate of the ground truth, namely, $$\mathbb{E}[\bm{M}]=\bm{M}^{\star},$$ where
the expectation is taken over the randomness in $\Omega$.


As a result, the proposed spectral method proceeds by computing the rank-$r$ SVD $\bm{U}\bm{\Sigma}\bm{V}^{\top}$ of the matrix $\bm{M}$ constructed in \eqref{eq:rescaled-M-matrix-completion-l2}, and employing $\bm{U}\in \mathbb{R}^{n_1\times r}$, $\bm{V}\in \mathbb{R}^{n_2\times r}$  and $\bm{U}\bm{\Sigma}\bm{V}^{\top}$ as estimates of $\bm{U}^{\star}$, $\bm{V}^{\star}$ and $\bm{M}^{\star}$, respectively.






\subsection{Performance guarantees}\label{sec:mc-l2-theory}

As before, whether the subspace $\bm{U}$ (resp.~$\bm{V}$) is close
to $\bm{U}^{\star}$ (resp.~$\bm{V}^{\star}$) relies crucially on
the size of the perturbation $\|\bm{M}-\bm{M}^{\star}\|$.
Therefore, we begin by developing an upper bound on this quantity; the proof is based on the matrix Bernstein inequality and is postponed to Section~\ref{sec:proof-auxiliary-lemmas-MC-L2}.

\begin{lemma}
\label{lemma:mc-noise-bound}
Consider the settings in Section~\ref{sec:problem-formulation-MC}. Suppose that $n_{2}p\geq C\mu r\log n_{2}$ for some
 constant $C>0$. Then with probability at least
	$1-O(n_{2}^{-10})$, the matrix $\bm{M}$ constructed in \eqref{eq:rescaled-M-matrix-completion-l2} obeys
%
\[
	\big\| \bm{M} -\bm{M}^{\star} \big\|
	\lesssim\sqrt{\frac{\mu r\log n_{2}}{n_{1}p}}\, \big\| \bm{M}^{\star} \big\|.
\]
%
\end{lemma}
%
With this perturbation bound in place, we are equipped to apply Wedin's $\sin\bm{\Theta}$ theorem to obtain the following results.  The condition on the sample size in Theorem~\ref{thm:mc-l2-subspace} is stronger than that in Lemma~\ref{lemma:mc-noise-bound}, as we need to control the eigengap in the following theorem.
%
\begin{theorem}
\label{thm:mc-l2-subspace}
Consider the settings in Section~\ref{sec:problem-formulation-MC}.
Suppose that $n_{1}p\geq C_{1}\kappa^{2}\mu r\log n_{2}$ for some sufficiently
large constant $C_{1}>0$. Then with probability exceeding $1-O(n_2^{-10})$,
%
\begin{align*}
\max\Big\{ \mathsf{dist}\left(\bm{U},\bm{U}^{\star}\right),\mathsf{dist}\left(\bm{V},\bm{V}^{\star}\right)\Big\}  & \lesssim\kappa\sqrt{\frac{\mu r \log n_{2}}{n_{1}p}} .
%;\\
%\max\Big\{ \mathsf{dist}_{\mathrm{F}}\left(\bm{U},\bm{U}^{\star}\right),\mathsf{dist}_{\mathrm{F}}\left(\bm{V},\bm{V}^{\star}\right)\Big\}  & \lesssim\kappa\sqrt{\frac{\mu r^{2} \log n_{2}}{n_{1}p}}.
\end{align*}
%
\end{theorem}
%
\begin{proof}
%	
As a direct consequence of Lemma \ref{lemma:mc-noise-bound}, one has
%
\[
	\big\| \bm{M} -\bm{M}^{\star} \big\|
	\lesssim\sqrt{\frac{\mu r\log n_{2}}{n_{1}p}} \,\big\| \bm{M}^{\star}\big\| \le \Big( 1 - \frac{1}{\sqrt{2}} \Big) \sigma_{r}(\bm{M}^{\star}),
\]
%
provided that $n_{1}p\geq C_{1}\kappa^{2}\mu r\log n_{2}$ for some  large enough constant $C_{1}>0$.
Apply Wedin's theorem (cf.~\eqref{eq:wedin-simpler-friendly}) and Lemma \ref{lemma:mc-noise-bound} to obtain
%
\begin{align*}
 & \max\Big\{ \mathsf{dist}\left(\bm{U},\bm{U}^{\star}\right),\mathsf{dist}\left(\bm{V},\bm{V}^{\star}\right) \Big\}
	\leq  \frac{2\big\|\bm{M} -\bm{M}^{\star}\big\|}{\sigma_{r}(\bm{M}^{\star})}
	%\\ % & \qquad\qquad
	%\lesssim\frac{\sqrt{\frac{\mu r\log n_{2}}{n_{1}p}}\,\big\|\bm{M}^{\star}\big\|}{\sigma_{r}(\bm{M}^{\star})}\asymp
	\lesssim \kappa\sqrt{\frac{\mu r\log n_{2}}{n_{1}p}}
\end{align*}
%
as claimed.
%The proof for the $\mathsf{dist}_{\mathrm{F}}(\cdot,\cdot)$ bounds follows from the same argument and is hence omitted.
%
\end{proof}


As an important implication of Theorem~\ref{thm:mc-l2-subspace}, once the sample size exceeds
$$pn_1n_2 \gg \kappa^{2}\mu r n_2  \log n_{2}, $$
then the spectral estimate achieves consistent estimation in the sense that
%
\[
	\max\Big\{ \mathsf{dist}\left(\bm{U},\bm{U}^{\star}\right),\mathsf{dist}\left(\bm{V},\bm{V}^{\star}\right)\Big\} = o(1).
\]
%
Given that $pn_1n_2\gtrsim \mu n_2 r \log n_2$ is an information-theoretic sampling requirement for reliable matrix completion when $r=o(n_1/\log n_2)$ \citep{candes2010NearOptimalMC},
Theorem~\ref{thm:mc-l2-subspace} confirms the near optimality of spectral methods---in terms of the scaling with $n_1$, $n_2$ and $p$---when it comes to consistent subspace estimation.



Before moving forward to the proof, we further characterize the statistical accuracy of $\bm{U}\bm{\Sigma}\bm{V}^{\top}$ in estimating the unknown matrix $\bm{M}^{\star}$. Accomplishing this only requires Lemma~\ref{lemma:mc-noise-bound}, without any need of the singular subspace perturbation theory. This result will also come in handy when we turn to discussing entrywise estimation accuracy in Chapter~\ref{cha:Linf-theory}.
%
\begin{theorem}
\label{thm:mc-l2-matrix}
Consider the settings in Section~\ref{sec:problem-formulation-MC}.
Suppose that $n_{2}p\geq C\mu r\log n_{2}$ for some sufficiently large constant
$C>0$. Then with probability at least $1-O(n_{2}^{-10})$, one has
%
\begin{align*}
%\|\bm{U}\bm{\Sigma}\bm{V}^{\top}-\bm{M}^{\star}\| & \lesssim\sqrt{\frac{\mu r}{n_{1}p}\log n_{2}}\,\sigma_{1}(\bm{M}^{\star});\\
\|\bm{U}\bm{\Sigma}\bm{V}^{\top}-\bm{M}^{\star}\|_{\mathrm{F}} & \lesssim\sqrt{\frac{\mu r^{2}\log n_{2}}{n_{1}p}}\, \big\| \bm{M}^{\star} \big\|.
\end{align*}
%
\end{theorem}
%
\begin{proof}
First, note
%
\begin{align*}
\|\bm{U}\bm{\Sigma}\bm{V}^{\top}-\bm{M}^{\star}\| &\leq\|\bm{U}\bm{\Sigma}\bm{V}^{\top}- \bm{M} \|+\|\bm{M} -\bm{M}^{\star}\| \leq2\|\bm{M}-\bm{M}^{\star}\|,
\end{align*}
%
where the first inequality comes from the triangle inequality, and the second inequality follows from the fact that $\bm{U}\bm{\Sigma}\bm{V}^{\top}$ is the best rank-$r$ approximation to $\bm{M}$,
i.e.,
$$\|\bm{U}\bm{\Sigma}\bm{V}^{\top}-\bm{M}\|=\min_{\bm{Z}: \mathsf{rank}(\bm{Z})\leq r}\|\bm{Z}-\bm{M}\| \leq \|\bm{M}-\bm{M}^{\star}\|. $$
Additionally, it is observed that $\bm{U}\bm{\Sigma}\bm{V}^{\top}-\bm{M}^{\star}$ has rank at most $2r$, which implies
%
\[
	\|\bm{U}\bm{\Sigma}\bm{V}^{\top}-\bm{M}^{\star}\|_{\mathrm{F}}\leq\sqrt{2r}\|\bm{U}\bm{\Sigma}\bm{V}^{\top}-\bm{M}^{\star}\| \leq 2\sqrt{2r} \big\|\bm{M} - \bm{M}^{\star} \big\|.
\]
%
This combined with Lemma~\ref{lemma:mc-noise-bound} immediately concludes the proof. \end{proof}

 


\subsection{Proof of auxiliary lemmas}
\label{sec:proof-auxiliary-lemmas-MC-L2}

\paragraph{Proof of Lemma~\ref{claim:mu-incoherence}.}
First of all, the  $\|\cdot\|_{2,\infty}$ norm of $\bm{M}^{\star}$
can be upper bounded by
%
\begin{equation*}
\|\bm{M}^{\star}\|_{2,\infty}=\|\bm{U}^{\star}\bm{\Sigma}^{\star}\bm{V}^{\star\top}\|_{2,\infty}\leq\|\bm{U}^{\star}\|_{2,\infty}\|\bm{\Sigma}^{\star}\| \, \|\bm{V}^{\star}\|\leq\sqrt{\frac{\mu r}{n_{1}}} \big\| \bm{M}^{\star} \big\|.
\end{equation*}
%
Here, the first inequality arises from the elementary bounds
$\|\bm{A}\bm{B}\|_{2,\infty}\leq\|\bm{A}\|_{2,\infty}\|\bm{B}\|$
and $\|\bm{A}\bm{B}\|\leq\|\bm{A}\|\, \|\bm{B}\|$, whereas the last relation
uses Definition~\ref{assump:mc-incoherence}, the orthonormality
of $\bm{V}^{\star}$, and identifies $\|\bm{\Sigma}^{\star}\|$ with
$\| \bm{M}^{\star} \|$. The  bound on $\|\bm{M}^{\star\top}\|_{2,\infty}$ can be derived analogously and is omitted for brevity.



In addition, the matrix $\bm{M}^{\star}$ is elementwise bounded by
%
\begin{equation*}
\|\bm{M}^{\star}\|_{\infty}=\|\bm{U}^{\star}\bm{\Sigma}^{\star}\bm{V}^{\star\top}\|_{\infty}\leq\|\bm{U}^{\star}\|_{2,\infty}\|\bm{V}^{\star}\|_{2,\infty} \|\bm{\Sigma}^{\star}\|
	\leq\frac{\mu r}{\sqrt{n_{1}n_{2}}} \big\| \bm{M}^{\star} \big\|.
\end{equation*}
%
Here, the first inequality follows from the fact $\|\bm{A}\bm{B}^{\top}\|_{\infty}\leq\|\bm{A}\|_{2,\infty}\|\bm{B}\|_{2,\infty}$
and the aforementioned one $\|\bm{A}\bm{B}\|_{2,\infty}\leq\|\bm{A}\|_{2,\infty}\|\bm{B}\|$,
while the last inequality again relies on Definition~\ref{assump:mc-incoherence}.
%\end{proof}


\paragraph{Proof of Lemma~\ref{lemma:mc-noise-bound}.} Note that the matrix $\bm{E}\coloneqq p^{-1}\mathcal{P}_{\Omega}(\bm{M}^{\star})-\bm{M}^{\star}$
can be expressed as the sum of $n_{1}n_{2}$ i.i.d.~random matrices
%
\[
  \frac{1}{p}  \mathcal{P}_{\Omega}(\bm{M}^{\star})-\bm{M}^{\star}=\sum_{i=1}^{n_{1}}\sum_{j=1}^{n_{2}}\underbrace{\big(p^{-1} \delta_{i,j} -1 \big)M_{i,j}^{\star}\bm{e}_{i}\bm{e}_{j}^{\top}}_{\eqqcolon\bm{X}_{i,j}}.
\]
%
Here, $\delta_{i,j}$ (which indicates whether the $(i,j)$-th entry is observed) follows an independent Bernoulli distribution
with parameter~$p$, and $\bm{e}_{i}$ stands for the $i$-th standard
basis vector of appropriate dimensions. It is easily seen that for each $(i,j)$,
%
\[
	\mathbb{E}[\bm{X}_{i,j}]=\bm{0}\quad\text{and}\quad\|\bm{X}_{i,j}\|\leq  \frac{1}{p} \|\bm{M}^{\star}\|_{\infty}\leq\frac{\mu r}{p\sqrt{n_{1}n_{2}}} \big\| \bm{M}^{\star} \big\|,
\]
%
where the last relation results from the entrywise upper bound (\ref{eq:mc-entry-upper})
on $\bm{M}^{\star}$. In order to apply the matrix Bernstein inequality (cf.~Corollary~\ref{thm:matrix-Bernstein-friendly}),
we need to control the variance statistic
%
\[
	v\coloneqq\max \Big\{ \Big\Vert \sum\nolimits _{i,j}\mathbb{E}\left[\bm{X}_{i,j}\bm{X}_{i,j}^{\top}\right]\Big\Vert ,\Big\Vert \sum\nolimits _{i,j}\mathbb{E}\left[\bm{X}_{i,j}^{\top}\bm{X}_{i,j}\right] \Big\Vert \Big\} .
\]
%
Regarding the first variance term, we have
%
\begin{align*}
\sum\nolimits _{i,j}\mathbb{E}\left[\bm{X}_{i,j}\bm{X}_{i,j}^{\top}\right]
	& =\sum\nolimits _{i,j}\mathbb{E}\left[ \big( p^{-1} \delta_{i,j} -1 \big)^{2}(M_{i,j}^{\star})^{2}\bm{e}_{i}\bm{e}_{j}^{\top}\bm{e}_{j}\bm{e}_{i}^{\top}\right]\\
 & =\frac{1-p}{p}\sum\nolimits _{i,j}\big(M_{i,j}^{\star}\big)^{2}\bm{e}_{i}\bm{e}_{i}^{\top}=\frac{1-p}{p}\sum_{i=1}^{n_{1}}\|\bm{M}_{i,\cdot}^{\star}\|_{2}^{2}\bm{e}_{i}\bm{e}_{i}^{\top} \\
	& \preceq \frac{1-p}{p}  \|\bm{M}^{\star}\|_{2,\infty}^{2} \bm{I}_{n_1} \preceq \frac{\mu r}{n_1p}  \|\bm{M}^{\star}\|^{2} \,\bm{I}_{n_1} .
\end{align*}
%
Here, the first identity arises from the definition of $\bm{X}_{i,j}$,
the second one calculates the variance of Bernoulli random variables,
and the last line relies on the upper bound~(\ref{eq:mc-row-norm-upper}).
Similarly, the second term in the variance statistic enjoys the following characterization:
%
\[
	\sum\nolimits _{i,j}\mathbb{E}\left[\bm{X}_{i,j}^{\top}\bm{X}_{i,j}\right] \preceq \frac{\mu r}{n_2p} \|\bm{M}^{\star}\|^{2} \,\bm{I}_{n_2}.
\]
%
Taking the above  relations together and recalling that $n_1\leq n_2$ give
%
\begin{align*}
v & %\leq\frac{1}{p}\max\left\{ \|\bm{M}^{\star}\|_{2,\infty}^{2},\|\bm{M}^{\star\top}\|_{2,\infty}^{2}\right\}
	% \leq\frac{1}{p}\left\{ \frac{\mu r}{n_{1}} \big\| \bm{M}^{\star} \big\|^2,  \frac{\mu r}{n_{2}} \big\| \bm{M}^{\star} \big\|^2 \right\} \\
	\leq\frac{\mu r}{n_{1}p} \big\| \bm{M}^{\star} \big\|^2.
\end{align*}
% 


With the above bounds in place, invoking matrix Bernstein (see Corollary~\ref{thm:matrix-Bernstein-friendly}) reveals that: with probability at least $1-O(n_2^{-10})$,
%
\begin{align*}
	\| \bm{E}\| & \lesssim \sqrt{\frac{\mu r \| \bm{M}^{\star} \|^2 \log n_{2}}{n_{1}p} }+\frac{\mu r \| \bm{M}^{\star} \| \log n_{2}}{p\sqrt{n_{1}n_{2}}}
	\asymp \sqrt{\frac{\mu r \| \bm{M}^{\star} \|^2 \log n_{2}}{n_{1}p} } ,
\end{align*}
%
where the last inequality is valid as long as  $n_{2}p \gtrsim \mu r\log n_{2}$.






\paragraph{Proof of Theorem~\ref{thm:CI-noisy-matrix-completion}.} 

Given that $\eta_{i,j}$ is a Gaussian random variable and hence possibly
unbounded, we find it convenient to introduce a truncated version
as follows
%
\[
	\widetilde{\eta}_{i,j}=\eta_{i,j}\mathbbm{1}\big\{|\eta_{i,j}|\leq 5\sigma_{\eta}\sqrt{\log n}\big\},\qquad1\leq i,j\leq n.
\]
%
and 
%
\[
	\widetilde{M}_{i,j} = \begin{cases} \frac{1}{p} \big( M_{i,j}^{\star} + \widetilde{\eta}_{i,j} \big), \qquad & \text{if } (i,j)\in \Omega, \\ 0, &\text{else}. \end{cases}
\]
%
Repeating the analysis in Section~\ref{sec:proof-eqn-noise-matrix-E-bound-denoising}, we can show that 
%
\[
	\mathbb{P}\big\{ \bm{M} = \widetilde{\bm{M}} \big\}
	= \mathbb{P} \big\{ \eta_{i,j} = \widetilde{\eta}_{i,j}, \forall i,j\in [n] \big\} 
	\geq 1- n^{-10},
\]
%
meaning that  $\bm{M}$ and $\widetilde{\bm{M}}$ are equivalent with high probability. 
As a result, we shall concentrate on validating the confidence interval computed based on $\widetilde{\bm{M}}$ in the subsequent analysis. %, unless otherwise noted 
Before proceeding, we record several key properties about  $\widetilde{\eta}_{i,j}$ as follows:
%
\begin{align}
	\mathbb{E}[ \widetilde{\eta}_{i,j}] = 0, ~~~ \mathbb{E}[ \widetilde{\eta}^2_{i,j}]=(1-o(1)) \sigma_{\eta}^2, 
	 ~~~ \big|\widetilde{\eta}_{i,j}\big|\leq  5\sigma_{\eta}\sqrt{\log n}.  
	 \label{eq:eta-property-CI-MC}
\end{align}
%
 


The proof follows by invoking Theorem~\ref{thm:CI-general}, as long as the conditions required therein are satisfied. 
To begin with, the associated variance parameters are given by
%
\begin{align*}
\sigma_{i,j}^{2} & \coloneqq\mathbb{E}\left[\big(\widetilde{M}_{i,j}-M_{i,j}^{\star}\big)^{2}\right]\\
 & =p\mathbb{E}\left[\left(\frac{1-p}{p}M_{i,j}^{\star}+\frac{1}{p}\widetilde{\eta}_{i,j}\right)^{2}\right]+(1-p)\big(M_{i,j}^{\star}\big)^{2}\\
 & = 
\frac{\left(1-p\right)^{2}}{p} \big(M_{i,j}^{\star}\big)^{2}+\frac{1}{p}\mathbb{E}\left[\widetilde{\eta}_{i,j}^{2}\right]+(1-p)\big(M_{i,j}^{\star}\big)^{2} \\
 & =\frac{1-p}{p}\big(M_{i,j}^{\star}\big)^{2}+\frac{1-o(1)}{p}\sigma_{\eta}^{2}, 
\end{align*}
%
thus leading to
%
\begin{align*}
\sigma_{\min}^{2} & \coloneqq \min_{i,j}\sigma_{i,j}^{2}=\frac{1-p}{p}\min_{i,j}\big(M_{i,j}^{\star}\big)^{2}+\frac{1-o(1)}{p}\sigma_{\eta}^{2},\\
\sigma^{2} & \coloneqq \max_{i,j}\sigma_{i,j}^{2}=\frac{1-p}{p} \|\bm{M}^{\star}\|_{\infty}^2 +\frac{1-o(1)}{p}\sigma_{\eta}^{2}.
\end{align*}
%
Apparently, $\sigma^2/\sigma_{\min}^{2}=O(1)$ holds true under the assumptions of Theorem~\ref{thm:CI-noisy-matrix-completion}. 
In addition, the random variables $\{\widetilde{M}_{i,j}-M_{i,j}^{\star}\}$ are all bounded obeying
%
\[
\big|\widetilde{M}_{i,j}-M_{i,j}^{\star}\big|\leq\frac{1-p}{p}\big|M_{i,j}^{\star}\big|+\frac{1}{p}\big|\widetilde{\eta}_{i,j}\big|
\leq\frac{(1-p)\|\bm{M}^{\star}\|_{\infty}+\sigma_{\eta}\sqrt{5\log n}}{p}\eqqcolon B.
\]
%
These bounds readily imply that
%
\begin{equation}
	\frac{B}{\sigma}\asymp\frac{~\frac{(1-p)\|\bm{M}^{\star}\|_{\infty}+\sigma_{\eta}\sqrt{5\log n}}{p}~}{\sqrt{\frac{1-p}{p}}\|\bm{M}^{\star}\|_{\infty}+\frac{1}{\sqrt{p}}\sigma_{\eta}}\lesssim\sqrt{\frac{\log n}{p}}. 
	\label{eq:B-sigma-ratio-noisy-MC}
\end{equation}
%



Moving to the condition $\sigma \kappa\sqrt{n\log n}\lesssim |\lambda_{r}^{\star}|$ in Theorem~\ref{thm:CI-noisy-matrix-completion}, it can be guaranteed if
	%
	\begin{equation*}
		\|\bm{M}^{\star}\|_{\infty}\kappa\sqrt{\frac{(1-p)n\log n}{p}}
		\lesssim|\lambda_{r}^{\star}| 
		\quad \text{and} \quad
		\sigma_{\eta}\kappa\sqrt{\frac{n\log n}{p}}\lesssim|\lambda_{r}^{\star}|.  
		\label{eq:Mstar-noise-requirement-noisy-MC-135}
	\end{equation*}
	%
Given the assumption  $\frac{\max_{k,l}|M_{k,l}^{\star}|}{\min_{k,l}|M_{k,l}^{\star}|}=O(1)$, one has 
%
\begin{equation}
	\label{eq:Mstar-inf-F-relation-noisy-MC}
	\|\bm{M}^{\star}\|_{\infty} \asymp \frac{1}{n} \|\bm{M}^{\star}\|_{\mathrm{F}} 
	\leq \frac{\sqrt{r}}{n} \|\bm{M}^{\star}\| 
	= \frac{\kappa \sqrt{r}}{n} |\lambda_{r}^{\star}|. 
\end{equation}
%
As a consequence, the condition $\sigma \kappa\sqrt{n\log n}\lesssim |\lambda_{r}^{\star}|$ can be ensured under Condition~\eqref{eq:Mstar-noise-requirement-noisy-MC}. 
%


It remains to certify Condition~\eqref{eq:Uj-Ui-lower-bound-thm}. 
By virtue of the above calculations of $\sigma$ and $B$ as well as the property \eqref{eq:B-sigma-ratio-noisy-MC}, it is easily seen that Condition~\eqref{eq:Uj-Ui-lower-bound-thm} is valid as long as the following holds:
%
\begin{align*}
\frac{\big\|\bm{U}_{j,\cdot}^{\star}\big\|_{2}^{2}+\big\|\bm{U}_{i,\cdot}^{\star}\big\|_{2}^{2}}{\big\|\bm{U}^{\star}\big\|_{\mathrm{F}}^{2}} & \gtrsim\frac{\kappa^{2}\mu^{2}r^{2}\log^{5/2}n}{n\sqrt{np}}+\frac{\sqrt{1-p}\,\|\bm{M}^{\star}\|_{\infty}\mu^{2}r\kappa^{3}\log^{3}n}{|\lambda_{r}^{\star}|\sqrt{np}}\\
 & \qquad+\frac{\sigma_{\eta}\mu^{2}r\kappa^{3}\log^{3}n}{|\lambda_{r}^{\star}|\sqrt{np}}.
\end{align*}
%
Taking this together with the relation \eqref{eq:Mstar-inf-F-relation-noisy-MC}, 
we can demonstrate straightforwardly that Condition~\eqref{eq:Uj-Ui-lower-bound-thm} is guaranteed to hold as long as Condition~\eqref{eq:Ui-Uj-lower-bound-thm-noisy-MC-146} is satisfied. 
This completes the proof. 




%\paragraph{Statistical guarantees.} 
